\newglossaryentry{adhoc}{
  name=ad hoc,
  description={Expresión para referirse a que se hace con un fin determinado}
}

\newglossaryentry{aaa}{
  name=AAA,
  description={Clasificación para videojuegos desarrollados con un alto presupuesto}
}

\newglossaryentry{astar}{
  name=A*,
  description={Algoritmo de búsqueda de caminos en grafo que hace uso de heurísticas}
}

\newglossaryentry{octree}{
  name=octree,
  description={Estructura de datos en árbol donde cada nodo interno tiene exactamente ocho hijos. Se utiliza principalmente para la partición jerárquica y eficiente del espacio tridimensional}
}

\newglossaryentry{llvm}{
  name=LLVM,
  description={Infraestructura de compilación modular diseñada para optimizar el tiempo de compilación, enlace y ejecución de programas. Su nombre no es un acrónimo}
}

\newglossaryentry{noholonomico}{
  name={no holonómico},
  description={Entidad cuyo movimiento está sujeto a restricciones cinemáticas y físicas, lo que le impide cambiar de dirección o velocidad de manera instantánea (por ejemplo, un vehículo con un radio de giro limitado)}
}

\newglossaryentry{stack}{
  name={stack},
  description={Región de memoria gestionada de forma automática y secuencial (último en entrar, primero en salir). Se utiliza para almacenar variables locales y el contexto de las llamadas a funciones durante la ejecución de un programa}
}

\newglossaryentry{heap}{
  name={heap},
  description={Región de memoria utilizada para la asignación dinámica. Los bloques de memoria se reservan y liberan en tiempo de ejecución, ya sea manualmente o mediante un recolector de basura, sin un orden establecido}
}

\newglossaryentry{fnv1a}{
  name={FNV-1a},
  description={Función hash no criptográfica diseñada por Fowler, Noll y Vo, conocida por su rapidez y alta tasa de dispersión. Se utiliza frecuentemente en tablas hash y aplicaciones donde el rendimiento es crítico, ya que procesa los datos de forma eficiente mediante operaciones simples a nivel de bits (XOR y multiplicación)}
}

\newglossaryentry{racecondition}{
  name={race condition},
  description={Anomalía que ocurre en la programación concurrente o multihilo cuando dos o más subprocesos intentan acceder y modificar un recurso compartido de forma simultánea. El estado final del sistema se vuelve impredecible, ya que depende del orden exacto e incontrolable en el que el procesador planifique la ejecución de las instrucciones}
}

\newglossaryentry{voxel}{
  name={vóxel},
  description={Acrónimo de \textit{volumetric pixel} (píxel volumétrico). Constituye la unidad cúbica mínima procesable dentro de una matriz o cuadrícula tridimensional. Es el equivalente tridimensional del píxel y se utiliza ampliamente para discretizar y representar la ocupación de espacios volumétricos}
}

\newglossaryentry{bake}{
  name={bake},
  description={Proceso de precomputación realizado durante la fase de desarrollo (en el editor). Consiste en calcular de forma anticipada información costosa —como la generación de mallas de navegación, estructuras espaciales o iluminación estática— y almacenarla en disco, evitando así realizar dichos cálculos durante el tiempo de ejecución del sistema}
}

\newglossaryentry{deadlock}{
  name={deadlock},
  description={Situación de interbloqueo temporal o permanente. En el contexto de la navegación de agentes, ocurre cuando dos o más entidades bloquean mutuamente sus respectivas trayectorias, impidiendo que cualquiera de ellas pueda avanzar hacia su destino al verse incapaces de resolver la evasión geométrica de forma local}
}
